\begin{abstract}
A cascade can beat another policy at the same fallback call rate and still lose to direct model calls at the same expected spend. Evaluating this distinction from repeated standalone responses is difficult because the pairing of first and fallback answers can change both accuracy and cost. We derive sharp bounds on a gate's excess score over a query-independent mixture of direct services at each replay's own spend. The calculation uses linear assignments and returns an attainable solution and a matching lower certificate. On 24 held-out AIME questions, a length gate from a 14B to a 32B model gained 4.19 accuracy points over call-matched random escalation under independent replay. At equal output-token prices it lost 4.17 points to direct service, with losses under every allowed pairing of the archived draws. Across 5,000 question resamples the call-to-spend reversal recurred in 82.1\%. Its independent-replay advantage became positive when the fallback token price exceeded 1.93 times the first-stage price. Two published cascade evaluations provide context. FrugalGPT retained an advantage at 21 of 22 distinct test operating points, 13 by more than one score point. For AutoMix, switching from mean incremental benefit per cost to a common expected budget changed one of five seed-averaged choices, by at most 0.14 points. These results identify which comparisons depend on spending assumptions and quantify the size of that dependence.
\end{abstract}
